% ========================================================
% BÀI TẬP TOÁN 5 - KẾT NỐI TRI THỨC VỚI CUỘC SỐNG
% Bài 36: Tỉ số. Tỉ số phần trăm
% Định dạng: BTTeX Part III - Trắc nghiệm điền số ngắn
% ========================================================

% Câu 1
\begin{ex}
  Trong lớp có $15$ bạn nam và $20$ bạn nữ. Tỉ số của số bạn nam và số bạn nữ khi rút gọn có mẫu số là bao nhiêu?
  \shortans{4}
  \loigiai{
    $\frac{15}{20} = \frac{3}{4}$. Mẫu số là 4.
  }
\end{ex}

% Câu 2
\begin{ex}
  Ở câu trên, tỉ số của số bạn nữ và số bạn nam khi rút gọn có tử số là bao nhiêu?
  \shortans{4}
  \loigiai{
    $\frac{20}{15} = \frac{4}{3}$. Tử số là 4.
  }
\end{ex}

% Câu 3
\begin{ex}
  Viết tỉ số $\frac{3}{4}$ dưới dạng tỉ số phần trăm là bao nhiêu phần trăm?
  \shortans{75}
  \loigiai{
    $\frac{3}{4} = 75\%$.
  }
\end{ex}

% Câu 4
\begin{ex}
  Viết tỉ số $\frac{1}{2}$ dưới dạng tỉ số phần trăm là bao nhiêu phần trăm?
  \shortans{50}
  \loigiai{
    $\frac{1}{2} = 50\%$.
  }
\end{ex}

% Câu 5
\begin{ex}
  Viết tỉ số $\frac{2}{5}$ dưới dạng tỉ số phần trăm là bao nhiêu phần trăm?
  \shortans{40}
  \loigiai{
    $\frac{2}{5} = 40\%$.
  }
\end{ex}

% Câu 6
\begin{ex}
  Viết tỉ số $\frac{7}{10}$ dưới dạng tỉ số phần trăm là bao nhiêu phần trăm?
  \shortans{70}
  \loigiai{
    $\frac{7}{10} = 70\%$.
  }
\end{ex}

% Câu 7
\begin{ex}
  Viết tỉ số $\frac{1}{4}$ dưới dạng tỉ số phần trăm là bao nhiêu phần trăm?
  \shortans{25}
  \loigiai{
    $\frac{1}{4} = 25\%$.
  }
\end{ex}

% Câu 8
\begin{ex}
  Viết số $1$ dưới dạng tỉ số phần trăm là bao nhiêu phần trăm?
  \shortans{100}
  \loigiai{
    $1 = 100\%$.
  }
\end{ex}

% Câu 9
\begin{ex}
  Một tổ sản xuất có kế hoạch làm $200$ sản phẩm nhưng đã làm được $240$ sản phẩm. Tổ đó đã thực hiện được bao nhiêu phần trăm kế hoạch?
  \shortans{120}
  \loigiai{
    $(240 : 200) \times 100\% = 120\%$.
  }
\end{ex}

% Câu 10
\begin{ex}
  Ở câu trên, tổ đó đã vượt mức kế hoạch bao nhiêu phần trăm?
  \shortans{20}
  \loigiai{
    $120\% - 100\% = 20\%$.
  }
\end{ex}

% Câu 11
\begin{ex}
  Một khu vườn có $80$ cây, trong đó có $20$ cây cam. Tỉ số phần trăm của số cây cam so với tổng số cây là bao nhiêu phần trăm?
  \shortans{25}
  \loigiai{
    $(20 : 80) \times 100\% = 25\%$.
  }
\end{ex}

% Câu 12
\begin{ex}
  Kiểm tra $40$ sản phẩm thì có $2$ phế phẩm. Tỉ số phần trăm phế phẩm là bao nhiêu phần trăm?
  \shortans{5}
  \loigiai{
    $(2 : 40) \times 100\% = 5\%$.
  }
\end{ex}

% Câu 13
\begin{ex}
  Ở câu trên, tỉ số phần trăm sản phẩm đạt chuẩn là bao nhiêu phần trăm?
  \shortans{95}
  \loigiai{
    $100\% - 5\% = 95\%$.
  }
\end{ex}

% Câu 14
\begin{ex}
  Số học sinh giỏi của lớp 5A chiếm $45\%$ số học sinh cả lớp. Hỏi số học sinh không đạt loại giỏi chiếm bao nhiêu phần trăm?
  \shortans{55}
  \loigiai{
    $100\% - 45\% = 55\%$.
  }
\end{ex}

% Câu 15
\begin{ex}
  Một dung dịch muối có $15\text{ g}$ muối hòa tan trong $85\text{ g}$ nước. Tỉ số phần trăm muối trong dung dịch là bao nhiêu phần trăm?
  \shortans{15}
  \loigiai{
    Khối lượng dung dịch là $15 + 85 = 100\text{ g}$. Tỉ số phần trăm là $15\%$.
  }
\end{ex}

% Câu 16
\begin{ex}
  Tính: $25\% + 35\%$. Kết quả là bao nhiêu phần trăm?
  \shortans{60}
  \loigiai{
    $25\% + 35\% = 60\%$.
  }
\end{ex}

% Câu 17
\begin{ex}
  Tính: $80\% - 35\%$. Kết quả là bao nhiêu phần trăm?
  \shortans{45}
  \loigiai{
    $80\% - 35\% = 45\%$.
  }
\end{ex}

% Câu 18
\begin{ex}
  Tính: $15\% \times 4$. Kết quả là bao nhiêu phần trăm?
  \shortans{60}
  \loigiai{
    $15\% \times 4 = 60\%$.
  }
\end{ex}

% Câu 19
\begin{ex}
  Tính: $75\% : 3$. Kết quả là bao nhiêu phần trăm?
  \shortans{25}
  \loigiai{
    $75\% : 3 = 25\%$.
  }
\end{ex}

% Câu 20
\begin{ex}
  Giá một chiếc áo là $200\,000$ đồng, được giảm giá $10\%$. Số tiền được giảm là bao nhiêu nghìn đồng?
  \shortans{20}
  \loigiai{
    $200 \times 10\% = 20$ nghìn đồng.
  }
\end{ex}

% Câu 21
\begin{ex}
  Ở câu trên, số tiền phải trả sau khi giảm giá là bao nhiêu nghìn đồng?
  \shortans{180}
  \loigiai{
    $200 - 20 = 180$ nghìn đồng.
  }
\end{ex}

% Câu 22
\begin{ex}
  Dân số một xã là $5\,000$ người, sau một năm tăng thêm $1\%$. Số dân tăng thêm sau một năm là bao nhiêu người?
  \shortans{50}
  \loigiai{
    $5\,000 \times 1\% = 50$ người.
  }
\end{ex}

% Câu 23
\begin{ex}
  Một người gửi tiết kiệm $10$ triệu đồng với lãi suất $6\%$ một năm. Sau một năm số tiền lãi nhận được là bao nhiêu trăm nghìn đồng?
  \shortans{6}
  \loigiai{
    $10 \times 6\% = 0{,}6$ triệu = 6 trăm nghìn đồng.
  }
\end{ex}

% Câu 24
\begin{ex}
  Số bi xanh bằng $\frac{2}{3}$ số bi đỏ. Tỉ số của số bi đỏ và tổng số bi có mẫu số là bao nhiêu?
  \shortans{5}
  \loigiai{
    Bi đỏ chiếm 3 phần, tổng số bi là $2 + 3 = 5$ phần. Tỉ số là $\frac{3}{5}$. Mẫu số là 5.
  }
\end{ex}

% Câu 25
\begin{ex}
  Lớp 5B có $32$ học sinh, trong đó có $24$ học sinh thích môn Toán. Tỉ số phần trăm học sinh thích môn Toán là bao nhiêu phần trăm?
  \shortans{75}
  \loigiai{
    $(24 : 32) \times 100\% = 75\%$.
  }
\end{ex}
